\documentclass[10pt, a4paper]{article}
\usepackage{preamble}

\title{Probability II \\
    \large Problem Class Notes}
\author{Luke Phillips}
\date{January 2026}

\begin{document}

\maketitle

\newpage

\tableofcontents

\newpage

\section{Problems Class \texorpdfstring{$1$}{}}

\begin{problem}
    Application of Boole's inequality.
    Let $A_n$,
    $n \geq 1$,
    be a sequence of events on a probability space such that $\P(A_n) = 0$ for all $n$.
    Show that $\displaystyle\P\left(\bigcup_{n = 1}^{\infty}A_n\right) = 0$ as well.

    Let $B_n$,
    $n \geq 1$,
    be a sequence of events such that $\P(B_n) = 1$ for all $n$.
    Show that
    \[
    \P\left(\bigcap_{n = 1}^{\infty}B_n\right) = 1
    \]
    as well.

    \begin{solution}
        Recall Boole's inequality:
        If $C_n$,
        $n \geq 1$,
        is a sequence of events then
        \[
        \P\left(\bigcup_{n = 1}^{\infty}C_n\right) \leq \sum_{n = 1}^{\infty}\P(C_n).
        \]
        We infer that
        \[
        0 \leq \P\left(\bigcup_{n = 1}^{\infty}A_n\right) \leq \sum_{n = 1}^{\infty}\P(A_n) = 0.
        \]
        Thus
        \[
        \P\left(\bigcup_{n = 1}^{\infty}A_n\right) = 0
        \]
        as required.

        Now let $C_n = B_n ^ c$.
        Then
        \[
        \P(C_n) = 1 - \P(B_n) = 0
        \]
        for all $n$.
        Therefore,
        $\P\left(\bigcup_{n = 1}^{\infty}C_n\right) = 0$ by the previous part.
        By De Morgan's law
        \[
        0 = \bigcup_{n = 1}^{\infty}C_n = \bigcup_{n = 1}^{\infty}B_n ^ c = \left(\bigcap_{n = 1}^{\infty}B_n\right) ^ c
        \]
        Therefore $\P\left(\left(\bigcap_{n = 1}^{\infty}B_n\right) ^ c\right) = 0$ so $\P\left(\bigcap_{n = 1}^{\infty}B_n\right) = 1$.
    \end{solution}
\end{problem}

\begin{problem}
    Let $X$ be a random variable such that
    \[
    \P(X \geq 0) = 1\qquad\text{and}\qquad\E(X) = 0.
    \]
    Show that $\P(X = 0) = 1$.

    \begin{solution}
        For $n \geq 1$ integer,
        consider the event $\{X \geq 1 / n\} = \{\omega \in \Omega\,:\,X(\omega) \geq 1 / n\}$.
        These events are increasing because if $X \geq 1 / n$ then $X \geq \frac{1}{n + 1}$ as well
        \big($\frac{1}{n + 1} < \frac{1}{n}$\big).
        Notice that $\bigcup_{n = 1}^{\infty}\{X \geq 1 / n\} = \{X > 0\}$.

        Indeed,
        if $X(\omega) > 0$ then there is $n \geq 1$ such that $X(\omega) \geq 1 / n$ and so that $\{X > 0\} \subseteq \bigcup_{n = 1}^{\infty}\{X \geq 1 / n\}$.

        Conversely,
        if $X(\omega) \geq 1 / n$ for some $n \geq 1$ then $X(\omega) > 0$,
        i.e.
        \[
        \bigcup_{n = 1}^{\infty}\{X \geq 1 / n\} \subseteq \{X > 0\}.
        \]
        We have $\P(X > 0) = \lim_{n \to \infty}\P(X \geq 1 / n)$.
        As $X \geq 0$,
        by Markov's inequality
        \[
        0 \leq \P(X \geq 1 / n) \leq n\E(X) = 0
        \]
        so $\P(X \geq 1 / n) = 0$ and so $\P(X > 0) = 0$ as well.
        Finally,
        \[
        1 = \P(X \geq 0) = \P(X = 0) + \P(X > 0) = \P(X = 0) + 0 \iff \P(X = 0) = 1
        \]
        $\{X = 0\}$,
        $\{X > 0\}$ are disjoint so we can separate at first step.
    \end{solution}
\end{problem}

\begin{problem}
    Let $X$ and $Y$ be independent with distribution
    (law)
    \begin{align*}
        \P(X = -1) &= \P(X = +1) = \frac{1}{2} \\
        \P(Y = -1) &= \P(Y = +1) = \frac{1}{2}.
    \end{align*}
    Define $Z = X\cdot Y$.
    Then each pair $(X, Y)$,
    $(X, Z)$ or $(Y, Z)$ is independent
    (i.e. $X, Y, Z$ are pairwise independent).
    However,
    the triple $(X, Y, Z)$ is not mutually independent.
    \[
    \P(X = a, Y = b, Z = c) \neq \P(X = a)\P(Y = b)\P(Z = c).
    \]

    \begin{solution}
        The pair $(X, Y)$ is independent by design.
        Consider $(X, Z)$.
        Note $Z \in \{-1, +1\}$.
        For $a, b \in \{-1, 1\}$ we have
        \begin{align*}
            \P(X = a, Z = b) &= \P(X = a, XY = b) \\
            &= \P(X = a, Y = ab) && \text{(as $a ^ 2 = 1$)} \\
            &= \P(X = a)\P(Y = ab) && \text{(as $X, Y$ independent)} \\
            &= \frac{1}{2}\cdot \frac{1}{2} = \frac{1}{4}
        \end{align*}
        \[
        \P(X = a) = \frac{1}{2}
        \]
        for $a \in \{-1, 1\}$.
        $\P(Z = 1) = \P(XY = 1) = \P(\{X = 1, Y = 1\})$ or $\{X = 1, Y = -1\}$,
        both events are disjoint so
        \[
        \P(Z = 1) = \P(X = 1, Y = 1) + \P(X = -1, Y = -1) = \frac{1}{4} + \frac{1}{4} = \frac{1}{2}.
        \]
        Similarly,
        $\P(Z = -1) = 1 - \P(Z = 1) = \frac{1}{2}$ as $Z \in \{-1, 1\}$.
        So $\P(Z = b) = \frac{1}{2}$ for each $b \in \{-1, 1\}$.
        Thus $\P(X = a, Z = b) = \frac{1}{4} = \P(X = a)\P(Z = b)$.
        Therefore $(X, Z)$ independent.
        A similar argument gives $(Y, Z)$ independent.
        However,
        \[
        \P(Z = 1, X = 1, Y = 1) = \P(X = 1, Y = 1) = \frac{1}{4}
        \]
        however $\P(X = 1)\P(Y = 1)\P(Z = 1) = \frac{1}{2}\cdot\frac{1}{2}\cdot\frac{1}{2} = \frac{1}{8}$ not equal so collection of three variables is not independent.
    \end{solution}
\end{problem}

\newpage

\section{Problems Class \texorpdfstring{$2$}{}}

\begin{problem}
    Let $X_n$ be a sequence of IID random variables with the $\Exp(1)$ distribution:
    \[
    \P(X_n > x) = e ^ {-x}\text{ for } x \geq 0.
    \]
    Show
    \[
    \limsup_{n \to \infty}\frac{X_n}{\log{n}} = 1\text{ with probability $1$}.
    \]
    \[
    \limsup_{n \to \infty}a_n \leq a\quad \forall k \geq 1, a_n \leq a + \frac{1}{k}\text{ eventually}
    \]
    \[
    \limsup_{n \to \infty}a_n \geq a\quad \forall k \geq 1, a_n \geq a - \frac{1}{k}\text{ infinitely often}.
    \]
    If both hold,
    then $\limsup_{n \to \infty}a_k = a$.

    \begin{solution}
        Define some events.
        For $n \geq 1$,
        $k \geq 1$,
        \begin{align*}
            A_{n, k} &= \{X_n > (1 + 1 / k)\log{n}\} & B_{n, k} &= \{X_n > (1 - 1 / k)\log{n}\} \\
            A_k &= \{A_{n, k}\text{ i.o in $n$}\} & B_k &= \{B_{n, k}\text{ i.o in $n$}\}\text{ for $k \geq 1$} \\
            A &= \bigcup_{k = 1}^{\infty}A_k & B 
            &= \bigcap_{k = 1}^{\infty}B_k.
        \end{align*}
        Calculate $\P$
        \begin{align*}
            \P(A_{n, k}) &= \P(X_n > (1 + 1 / k)\log{n}) = e ^ {-(1 + 1 / k)\log{n}} = n ^ {-(1 + 1 / k)} \\
            \P(B_{n, k}) &=  \P(X_n > (1 - 1 / k)\log{n}) = e ^ {-(1 - 1 / k)\log{n}} = n ^ {-1 + 1 / k}.
        \end{align*}
        \begin{align*}
            \sum_{n = 1}^{\infty}\P(A_{n, k}) &= \sum_{n = 1}^{\infty}n ^ {-(1 + 1 / k)} \\
            \sum_{n = 1}^{\infty}\P(B_{n, k}) &= \sum_{n = 1}^{\infty}n ^ {-1 + 1 / k} = +\infty.
        \end{align*}
        The first Borel-Cantelli lemma implies $\P(A_{k}) = \P(A_{n, k}\text{ i.o in $n$}) = 0$.
        The second Borel-Cantelli lemma implies $\P(B_k) = \P(B_{n, k}\text{ i.o in $n$}) = 1$.
        \[
        \P(A) = \P\left(\bigcup_{k = 1}^{\infty}A_k\right) = 0
        \]
        by Boole's inequality
        \[
        \P(B) = \P\left(\bigcap_{k = 1}^{\infty}\right) = 1.
        \]
        Let $C = A ^ c \cap B$
        $\P(A ^ c) = 1$,
        $\P(B) = 1$ imply $\P(C) = 1$.
        De Morgan's laws:
        $A ^ c = \bigcap_{k = 1}^{\infty}A_k ^ c$
        \[
        A_k ^ c = \{A_{n, k}\text{ i.o in $n$}\} ^ c = \{A_{n, k} ^ c \text{ eventually in $n$}\}.
        \]
        \[
        A_{n, k}^{c} = \{X_n \leq (1 + 1 / k)\log{n}\}.
        \]

        Suppose $\omega \in C \implies \omega \in A ^ c$ and $\omega \in B$ so $\omega \in A_{k} ^ c$ for every $k$ so $\omega \in \{A_{k} ^ c\text{ eventually in $n$}$ for every $k \geq 1$ which is equivalent to $X_n(\omega) \leq (1 + 1 / k)\log{n}$ eventually in $n$ so $\frac{X_n(\omega)}{\log{n}} \leq 1 + 1 / k$ eventually in $n$.
        $\omega \in \bigcap_{k = 1}^{\infty}A_k ^ c$ implies $\limsup_{n \to \infty}\frac{X_n(\omega)}{\log{n}} \leq 1$.

        If $\omega \in B$ then $\omega \in B_k$ for every $k$ so $\omega \in \{B_{n, k}\text{ i.o}\}$ for every $k$ so $\{X_{n}(\omega) > (1 - 1 / k)\log{n}\text{ i.o}\}$ in $n$ for every $k$ then
        \[
        \limsup_{n \to \infty}\frac{X_n(\omega)}{\log{n}} \geq 1.
        \]
        $\omega \in C = A ^ c \cap B$ then $\limsup_{n \to \infty}\frac{X_n(\omega)}{\log{n}} \leq 1 \implies \limsup_{n \to \infty}\frac{X_n(\omega)}{\log{n}} = 1$ $\P(C) = 1$ which implies $\P\left(\limsup_{n \to \infty}\frac{X_n}{\log{n}} = 1\right) = 1$.
    \end{solution}
\end{problem}

\begin{problem}
    A group of $N$ students are comparing their birthdays.
    Find the expected number of students with a common birthday.

    \begin{solution}
        Let $X = \#\text{ students that don't share any birthdays}$.
        Want $N - X$,
        $\E(N - X) = N - \E(X)$.
        \begin{align*}
            X &= \sum_{k = 1}^{N}\mathbbm{1}_{\{\text{student $k$ does not share a birthday}\}}  \\
            \implies \E(X) &= \sum_{k = 1}^{N}\P(\text{student $k$ does not share a birthday}) \\
            &= N\P(\text{student $1$ doesn't share a birthday})
        \end{align*}
        \begin{align*}
            \P(\text{student $1$ has a different birthday to students $2, \dotsc, n$}) &= \P(\text{student $1$ and $2$ have different birthdays}) \\
            &= \left(\frac{364}{365}\right) ^ {N - 1} \\
            \E(X) &= N\left(\frac{364}{365}\right) ^ {N - 1} \implies \text{Want } N\left(1 - \left(\frac{364}{365}\right) ^ {N - 1}\right).
        \end{align*}
    \end{solution}
\end{problem}

\newpage

\section{Problems Class \texorpdfstring{$3$}{}}

\begin{problem}
    Suppose $X_n \to X$ in probability and $Y_n \to Y$ in probability.
    Show that $X_n + Y_n \to X + Y$ in probability.

    \begin{solution}
        We know $\P(|X_n - X| > \varepsilon) \to 0$ for every $\varepsilon > 0$.
        $\P(|Y_n - Y| > \varepsilon) \to 0$ for every $\varepsilon > 0$.
        Want to show $\P(|(X_n + Y_n) - (X + Y)| > \varepsilon) \to 0$ for every $\varepsilon > 0$.
        Recall the triangle inequality for numbers:
        \[
        |a + b| \leq |a| + |b|
        \]
        where $a, b \in \R$.
        \[
        |(X_n + Y_n) - (X + Y)| = |(X_n - X) + (Y_n - Y)| \leq |X_n - X| + |Y_n - Y|.
        \]
        Suppose $|(X_n + Y_n) - (X + Y)| > \varepsilon$.
        Then $|X_n - X| > \frac{\varepsilon}{2}$ or $|Y_n - Y| > \frac{\varepsilon}{2}$.
        The event $\{|(X_n + Y_n) - (X + Y)| > \varepsilon\}$ implies $\{|X_n - X| > \varepsilon / 2\} \cup \{|Y_n - Y| > \varepsilon / 2\}$.
        \begin{align*}
            \P(|(X_n + Y_n) - (X + Y)| > \varepsilon) &\leq \P(|X_n - X| > \varepsilon / 2\text{ or } |Y_n - Y| > \varepsilon / 2) \\
            &\leq \P(|X_n - X| > \varepsilon / 2) + \P(|Y_n - Y| > \varepsilon / 2) \\
            \lim_{n \to \infty}\P(|(X_n + Y_n) - (X + Y)|) &\leq \lim_{n \to \infty}\P(|X_n - X| > \varepsilon / 2) + \P(|Y_n - Y| > \varepsilon / 2)
        \end{align*}
        since both sides go to $0$ in probability.
    \end{solution}
\end{problem}

\begin{problem}
    Suppose $X_n \to X$ in probability and $X_n \to Y$ in probability.
    Show that $\P(X = Y) = 1$
    (limits in probability are unique).

    \begin{solution}
        Know $\P(|X_n - X| > \varepsilon) \to 0$ for all $\varepsilon > 0$,
        $\P(|X_n - Y| > \varepsilon) \to 0$ for all $\varepsilon > 0$.

        Consider the event $\{|X - Y| > \varepsilon\}$.
        \begin{align*}
            |X - Y| &= |(X - X_n) + (X_n - Y)| \\
            &\leq |X - X_n| + |X_n - Y| \\
            &= |X_n - X| + |X_n - Y|.
        \end{align*}
        If $|X - Y| > \varepsilon$ then $|X_n - X| > \varepsilon / 2$ or $|X_n - Y| > \varepsilon / 2$.
        \begin{align*}
            \P(|X - Y| > \varepsilon) &\leq \P(|X_n - X| > \varepsilon / 2\text{ or }|X_n - Y| > \varepsilon / 2) \\
            &\leq \P(|X_n - X| > \varepsilon / 2) + \P(|X_n - Y| > \varepsilon / 2) \\
            \P(|X - Y| > \varepsilon) &\leq \lim_{n \to \infty}\P(|X_n - X| > \varepsilon / 2) + \lim_{n \to \infty}\P(|X_n - Y| > \varepsilon / 2) = 0 + 0 \\
            \P(|X - Y| > \varepsilon) &= 0\quad \forall \varepsilon > 0.
        \end{align*}

        For integer $k \geq 1$,
        let
        \[
        A_k = \{|X - Y| > 1 / k\}.
        \]
        Found $\P(A_k) = 0$ for every $k$.
        \[
        A = \bigcup_{k = 1}^{\infty}A_k
        \]
        claim $A = \{|X - Y| > 0\}$.
        \[
        \P(A) = \P\left(\bigcup_{k = 1}^{\infty}A_k\right) \leq \sum_{k = 1}^{\infty}\P(A_k) = 0.
        \]
        so $\P(|X - Y| > 0) = 0$.
        Then $\P(|X - Y| = 0) = 1 - \P(|X - Y| > 0) = 1$.
    \end{solution}
\end{problem}

\begin{problem}
    Suppose $X_n \to X$ in $L ^ r$ and $Y_n \to Y$ in $L ^ r$ for $r \geq 1$.
    Show that $X_n + Y_n \to X + Y$ in $L ^ r$.

    \begin{solution}
        Want $\E(|(X_n + Y_n) - (X + Y)| ^ r) \to 0$.
        Know $\E(|X_n - X| ^ r) \to 0$,
        $\E(|Y_n - Y| ^ r) \to 0$.

        \begin{align*}
            |(X_n + Y_n) - (X + Y)| &= |(X_n - X) + (Y_n - Y)| \\
            &\leq |X_n - X| + |Y_n - Y|.
        \end{align*}
        Fact:
        suppose $a, b > 0$ and $r \geq 1$.
        Then
        \[
        (a + b) ^ r \leq 2 ^ r(a ^ r + b ^ r).
        \]
        Apply this fact to $a = |X_n - X|$ and $b = |Y_n - Y|$.
        Then
        \begin{align*}
            |(X_n + Y_n) - (X + Y)| ^ r &\leq (|X_n - X| + |Y_n - Y|) ^ r \\
            &\leq 2 ^ r(|X_n - X| ^ r + |Y_n - Y| ^ r) \\
            \E(|(X_n + Y_n) - (X + Y)| ^ r) &\leq 2 ^ r(\E(|X_n - X| ^ r) + \E(|Y_n - Y| ^ r)) \\
            \lim_{n \to \infty}\E(|(X_n + Y_n) - (X + Y)| ^ r) &\leq 2 ^ r\left(\lim_{n \to \infty}\E(|X_n - X| ^ r) + \lim_{n \to \infty}\E(|Y_n - Y| ^ r)\right) \\
            &= 2 ^ r \cdot 0 = 0
        \end{align*}
        which implies $X_n + Y_n \to X + Y$ in $L ^ r$.

        Example:
        Let $U \sim U(0, 1)$,
        $X_n = \sqrt{n}\mathbbm{1}_{\{U \leq 1 / n\}}$,
        $Y_n = \sqrt{n}\mathbbm{1}_{\{U \leq 1 / (2n)\}}$,
        $X_n \to 0$ almost surely and $Y_n \to 0$ almost surely.
        $L ^ 1$ converges
        \[
        \E(|X_n|) = \E(\sqrt{n}\mathbbm{1}_{\{U \leq 1 / n\}} = \sqrt{n}\P(U \leq 1 / n) = \frac{\sqrt{n}}{n} = \frac{1}{\sqrt{n}} \to 0
        \]
        \[
        \E(|Y_n|) = \frac{\sqrt{n}}{2n} = \frac{1}{2\sqrt{n}} \to 0
        \]
        so $X_n, Y_n \to 0$ in $L ^ 1$.
        \[
        X_nY_n = n\mathbbm{1}_{\{U \leq 1 / n\}}\mathbbm{1}_{\{U \leq 1 / (2n)\}} = n\mathbbm{1}_{\{U \leq 1 / (2n)\}}
        \]
        then $\E(|X_nY_n|) = n\P(U \leq 1 / (2n)) = \frac{n}{2n} = \frac{1}{2} \centernot\to 0$.
    \end{solution}
\end{problem}


\end{document}